\section{产品起源与技术栈}

\subsection{发展历程}

\begin{table}[H]
\centering
\caption{Claude Code 与 Codex 发展时间线}
\begin{tabular}{lp{5.5cm}p{5.5cm}}
\toprule
\textbf{时间节点} & \textbf{Claude Code} & \textbf{Codex} \\
\midrule
起源 & Anthropic 内部 @bcherny 的 side project，终端原型可调用 Claude API、读取文件、执行 bash 命令 & 最初的 Codex 是 12B GPT-3 微调模型，驱动了初代 GitHub Copilot \\
\addlinespace
公开发布 & 2025-02-24，Research Preview，使用 Claude 3.7 Sonnet & 2025-04-16，Codex CLI 发布 \\
\addlinespace
最新版本 & Opus 4.6（1M token 上下文） & GPT-5.3-Codex（2026-02-05），OpenAI 称其为"第一个参与了自身创建的模型" \\
\bottomrule
\end{tabular}
\end{table}

\begin{knowledgebox}{内部采用速度}
Claude Code 原型在 Anthropic 内部发布后，第五天就有一半团队在使用。这种自发的内部采用往往是产品 product-market fit 的强信号。
\end{knowledgebox}

\subsection{技术栈对比}

\begin{table}[H]
\centering
\caption{技术栈对比}
\begin{tabular}{lcc}
\toprule
\textbf{维度} & \textbf{Claude Code} & \textbf{Codex CLI} \\
\midrule
编程语言 & TypeScript & Rust \\
UI 框架 & React + Ink（终端 UI） & Ratatui（Rust TUI 库） \\
分发方式 & 单个 Bun 可执行文件 & 原生二进制 \\
上下文窗口 & 1M tokens & -- \\
\bottomrule
\end{tabular}
\end{table}

作者注意到 Claude Code CLI 偶有小 glitch，但不影响实际编码体验。两者本质上都是围绕底层模型 API 的轻量包装（thin wrapper）。

\begin{knowledgebox}{Anthropic 收购 Bun}
Anthropic 于 2025 年 12 月收购了 JavaScript 运行时 Bun，Claude Code 因此可以打包为单个 Bun 可执行文件分发，简化了安装流程。
\end{knowledgebox}

\subsection{本章小结}

Claude Code 源于内部工具的自然生长，技术栈偏向 TypeScript 生态；Codex CLI 选择了 Rust 以追求性能和可移植性。两者都是模型 API 的轻量封装，产品差异更多体现在模型能力和交互设计上。

